\answer{ex:15:1}{(a)~$2\cdot10^5$ por neurona, $6\cdot10^{12}$ para el circuito. (b)~$\approx8$~bit/s por cable, como mínimo $2.5\cdot10^4$~s $\approx7$~h.}
\answer{ex:15:2}{Factor $2$: correlación entre vecinos $0.943$, $\lambda$ de $4.18$ a $7.3\cdot10^{-4}$, cociente $\approx5.7\cdot10^3$; factor $4$: correlación $0.8$, cociente $\approx94$.}
\answer{ex:15:3}{(a)~$\lambda=0.988$ y $0.808$: $82$ y $4.7$ movimientos. (b)~$x_s^*=0.690$, $x_f^*=0.172$, suma $0.862$.}
\answer{ex:15:4}{(a)~$116$ movimientos. (b)~$19$. Un modelo de un solo proceso con suma nula no da ningún ahorro.}
\answer{ex:15:5}{(a)~$G=50$~s$^{-1}$, frente al límite de $10.5$~s$^{-1}$ sin modelo. (b)~No: con $G=50$~s$^{-1}$ el retardo crítico es $d_c\approx103\ms$; con $d=150\ms$ hace falta $\hat k>0.445k$ (para cualesquiera $G$ y $d$, $\hat k>k/2$).}
\answer{ex:15:6}{El error baja del $0.5\%$ en $6$--$30$~s, con un piso de $\approx(6$--$9)\cdot10^{-4}$ frente a $7.5\cdot10^{-4}$ con los pesos óptimos; con $\eta=50$ los pesos aún difieren de los óptimos hasta en $0.2$ a lo largo de las direcciones mal condicionadas de $C$ (ejercicio~\ref{ex:15:2}).}
\answer{ex:15:7}{$\mathbf B=A\mathbf K$; $q_f/R\approx0.035$, $q_s/R\approx8.6\cdot10^{-4}$; con $4q_s$, $B_f\approx0.088$, $B_s\approx0.047$: el proceso lento aprende más del doble de rápido, y el rápido, más despacio.}
